% !TeX root = ../colloquium-2.tex

\phantomsection\label{sec_Proofs04}
\begin{proof}[Доказательство равносильности дифференцируемости и наличия производной. Вывод производной $f(x) = e^x$ и $f(x) = \cos x$. Доказательство непрерывности дифференцируемой функции. Доказательство того, что функция $f(x) = |x|$ не является дифференцируемой в нуле.]
    \subsection*{Доказательство равносильности дифференцируемости и наличия производной.}
    \begin{note}[amber]
        Функция $f$ дифференцируема в точке $a$ тогда и только тогда, когда в точке $a$ существует производная этой функции $f'(a)$. При этом $df(h)\vert_{x=a} = f'(a)h$. \\
        \begin{note}[purple]
            НЕОБХОДИМОСТЬ.\\
            Если функция $f$ дифференцируема, то справедливо равенство
            \[
            f(a+h)-f(a)=Ah+\alpha(h)h, \ \lim_{h\to 0}\alpha(h)=0.
            \]
            Разделив обе части этого равенства на $h$, получим
            \[
            \frac{f(a+h)-f(a)}{h}=A+\alpha(h).
            \]
            Переходя к пределу при $h\to 0$ и учитывая, что $\lim\limits_{h\to 0}\alpha(h)=0$, получим, что $f'(a)=A$.\\
            ДОСТАТОЧНОСТЬ.\\
            Если функция $f$ имеет производную, то есть существует предел
            \[
            \lim_{h\to 0}\frac{f(a+h)-f(a)}{h}=f'(a),
            \]
            то по определению предела справедливо равенство
            \[
            \frac{f(a+h)-f(a)}{h}=f'(a)+\alpha(h),
            \]
            где $\lim\limits_{h\to 0}\alpha(h)=0$. Домножив обе части этого равенства на $h$, мы получим равенство из определения дифференцируемой функции.
        \end{note}
    \end{note}

    \subsection*{Вывод производной $f(x) = e^x$ и $f(x) = \cos x$.}
    \begin{note}[amber]
        \begin{note}[purple]
            1) Пусть $f(x)=e^x$. Найдём $\dfrac{df}{dx}(a)$. Имеем по определению:
            \[
            \frac{df}{dx}(a)=\lim_{h\to 0}\frac{e^{a+h}-e^a}{h}
            = e^a \lim_{h\to 0}\frac{e^h-1}{h}
            = e^a.
            \]
            Так как это верно в любой точке области определения, то при всех $x \in \mathbb{R}$
            \[
            (e^x)' = e^x.
            \]
            2) Пусть $f(x)=\cos x$. Тогда
            \[
            f'(x)=\lim_{h\to 0}\frac{\cos(x+h)-\cos x}{h}
            = \lim_{h\to 0}\frac{-2\sin\frac{h}{2}\sin\frac{2x+h}{2}}{h}
            = - \lim_{h\to 0}\frac{\sin\frac{h}{2} \cdot \sin\frac{2x+h}{2}}{\frac{h}{2}}
            = -\sin x.
            \]
        \end{note}
    \end{note}

    \subsection*{Доказательство непрерывности дифференцируемой функции.}
    \begin{note}[amber]
        Пусть функция $f$ дифференцируема в точке $a$. Тогда $f$ непрерывна в точке $a$.
        \begin{note}[purple]
            Из определения дифференцируемости следует, что
            \[
            f(x)=f(a)+f'(a)(x-a)+o(x-a)\to f(a), \quad x\to a,
            \]
            то есть предел функции в точке $a$ равен её значению в этой точке, что по определению означает непрерывность.
        \end{note}
    \end{note}

    \subsection*{Доказательство того, что функция $f(x) = |x|$ не является дифференцируемой в нуле.}
    \begin{note}[amber]
        \begin{note}[purple]
            Действительно,
            \[
            \lim_{h\to 0^-}\frac{|h|}{h}=\lim_{h\to 0^-}\frac{-h}{h}=-1,
            \qquad
            \lim_{h\to 0^+}\frac{|h|}{h}=\lim_{h\to 0^+}\frac{h}{h}=1.
            \]
            Если бы производная в нуле существовала, то есть если бы существовал предел, то это было бы равносильно существованию и равенству односторонних пределов. 
            Таким образом, у функции $x=|x|$ не существует производной при $x=0$, что равносильно тому, что эта функция недифференцируема в нуле.
        \end{note}
    \end{note}
\end{proof}
